\section{We Antimoderns}

\begin{quotex}
In 1930, \textbf{Julius Evola} launched the periodical \emph{La Torre} which was subsequently shut down by the Fascist government of Italy. This was just after the publication of the Italian edition of \emph{Pagan Imperialism}, but before Revolt against the modern World. This essay, \textit{Noi Antimoderni}\footnote{\url{http://www.juliusevola.it/documenti/template.asp?cod=527}} is from the inaugural issue. While there is nothing unexpected here, it sets the scene for his subsequent works. The translation will appear in two parts.

Commentaries are featured with small letter.
\end{quotex}

In many ways today, a dark threat, understood in an always more precise sense, looms over the entire civilization of the West. In the crisis, without insisting on this or that special form, but rather on the structure of the entire modern world, it seems to foretell the symptoms of the end of a world, the decline of a culture.

While \textbf{Rene Guenon}, analyzes the imbalance and discomfort characteristic of this era, he, in fact, shows how its characteristics are precisely those of the Dark or Iron Age, foretold by ancient traditions.

\begin{quotex}
There are two important points to keep in mind:

\end{quotex}
\begin{enumerate}
\item \begin{quotex}
The so-called crisis is entwined in the very structure of the modern world. It is not a question of fixing a limited problem, winning an election, or purging the West of some group or form of thought. 
\end{quotex}
\item \begin{quotex}
Its characteristics can be known in advance by those who understand cosmic cycles. Everything that happens does so for a reason, by following certain principles, or rather, by the denial and neglect of real principles. 
\end{quotex}
\end{enumerate}
\textbf{Spengler} indicates how today there is in existence a fatally inflexible law, for which, like any organism, each civilization, after its rise and prospering, its fall and its petrification into a lifeless barbaric grandeur follow. 

After \textbf{Netzsche}, \textbf{Keyserling}, \textbf{Kalergi} indict the immoralism and the irrealism of the European soul, while \textbf{Benda} notes the treason of the intellectuals, the enslavement of the classes that were the legacy of a spiritual tradition, due to passion and political hatred. 

In fact, ancient certitudes totter everywhere, principles are uncertain everywhere, traditions are lost, spirits are divided, and dark, uncontrollable, and irrational forces impel and overwhelm men and society, playing with them through ideas, interests and passions that they delude themselves into pursuing.

\begin{quotex}
Haven't we had enough kvetching and jeremiads? We know things are awry, so what is the point of indulging in criticisms of this injustice or that stupid policy? There is a time for righteous anger and a time for sober reflection. Pace this critique, there must be men who stand above them.

\begin{itemize}
\item They know the ancient certitudes. 
\item They are certain about their principles. 
\item They seek to recover their traditions. 
\item Their spirits are unified in solidarity each other and continuity with their ancestors. 
\item They are aware of the dark and chaotic forces within themselves and learn to master them. 
\end{itemize}

Those who insist on simple or one-dimensional answers will get lost here.

\end{quotex}
That civilization, of which the Modern man was so proud, and in the name of which he believed the \emph{myth of progress} and had marched to the conquest of the world, that civilization is now faced with a \emph{reductio ad absurdum}, and with a reversal of the values which it had claimed for itself. 

Having thrown itself into the conquest of matter, it served its purpose only at the cost of materializing the spirit, of excluding any higher form of life, of amalgamating individuals in the tyranny of collective bodies, we would almost say \emph{subhumans} in their lack of face, rationality, and light, in their subjection to energies that from time to time, as though galvanizing with a temporary life fearful of dead or automated bodies. It hurls them against each other. 

The Christian attempt to give the West a religious tradition can only be considered a failure. Nostalgia, as shown in spirits like \textbf{Maritain}, \textbf{Guenon}, or \textbf{Berdyaev}, turn to the feudal and Catholic Middle Ages, putting aside perhaps the insurmountable distance between then and now, when Europe really had begun to organize itself under the two great symbols of Action and Contemplation. What does it matter that Christianity (without realizing it) served as a vehicle for the transmission of a Transcendent Wisdom, \emph{prior to every age}, and that the Church in its Rites, Symbols and Dogma conserved its deposit, if for some time now, there has been no consciousness corresponding to it? 

\begin{quotationx}
Evola's criticism is misguided here. There is no question of returning to the material conditions of a bygone age, but rather of recovering its essence. Maritain became infected with modernist ideas. Berdyaev, although not a Traditionalist per se, was nevertheless a man of Tradition. He called for a new Middle Ages. Guenon had a deeper understanding of Tradition and was beyond certain prejudices, as he adhered to a strict intellectuality.

It matters quite a bit if Christendom was the bearer of a Transcendent Wisdom. At the very least, it provides us with an understanding of Tradition that is ``close to home" in space and time. The exercise of extracting the Traditional elements from the thought, rites, symbols, and dogmas is very valuable and less prone to misunderstanding than the abstract study of Asian traditions. Conversely, we are dubious about the claim of any man to be a Traditionalist if he cannot recognize Tradition in those elements.

We agree that for quite some time, there has been no consciousness of such a Tradition in the visible representatives of Christianity today. \emph{A fortiori}, the Middle Ages are an embarrassment to them, and the pagan and Traditional elements in it were explicitly rejected by the Reformation. Hence, those of a nostalgic bent are disappointed and are forced to hold many mental reservations.

However, Evola's assertion that the Middle Ages were \emph{never} conscious of the Traditional element is simply absurd. A Tradition, by its very nature, requires a conscious elite to maintain it, and when that is lost, the Tradition is lost. It cannot run mechanically or automatically. Thus, Guenon's explanation that Christianity was originally an esoteric movement makes more sense. He points to The Divine Comedy as an initiatory poem, as well as other writings from the period. This is not to defend Christianity, but rather the notion of Tradition itself, as well as the intelligence and integrity of our European ancestors. 

\end{quotationx}
Christianity today is no longer of value to people than as a small faith and morality that everyone professes and everyone betrays; it is mediocre and bourgeois in Catholicism, weakened and stimulating practical achievements in the social rigour in Protestantism? It is only this regard that those who speak of \emph{Tradition} and of the \emph{Return to Tradition}, in fact, know even less of those who deny what Tradition is.

\begin{quotex}
Here Evola goes for the low lying fruit, repeating what has been said by others ad nauseum. At least he only takes a paragraph to do so. What is the value of criticizing the worst elements rather than the appreciating the best? In Ancient Greece and Rome, the bulk of the population consisted of serfs, slaves, and outcastes. Does he judge paganism by their superstitions or by the few priest-kings and warriors? Always remember there are different castes and types of men, who embody a Tradition to a greater or lesser degree. We should always look to the greatest exponents.

\end{quotex}
\textbf{Henri Massis}, who brought up the symbol of the Defence of the West and raised the alarm against the asianization of the Latin world, in reality, does not know what the East is nor what could be of value in the West as a principle of reintegration. He does not know how much of what he denies is in what he affirms, nor how much of what he affirms is in what he denies. We then are silent about everything that for some time we proclaimed concerning traditions and traditionalism, whether on this or that basis, those who exalting a Vatican Rome, others a Masonic Rome, or a Mazzinian and Giobertian Rome, who, raising strange taboos to the right and the left, launching attacks in the void, preparing with wordiness and the most implausible mess. Here, as elsewhere, the confusion of tongues is complete; the power of diagrams, formulas, and words, that, as entities created by Magic, no longer depending on those who created them, is almost without limit. 

Neither is sufficient. An amorphous need to escape to the Ahrimanic embrace of materialism, no longer meeting those supports that were given by the surviving traditions only on the assumption of interior and living relationships, has created an even more dangerous perversion in the unbalanced Western soul: \emph{neo-spiritualism}. 

From the various revivals of a suspect mysticism to the introduction of always more counterfeit exotic doctrines; from the newest spiritualistic superstition to the morbid interest in the problems and complications of the unconscious and psychoanalysis; from Intuitionism and Surrealism to the various messianic forms and to pseudo-religious and pseudo-occultists of the 1700s that swarmed at the margins of Protestantism; from humanitarian and universalistic ideologies to those of a \textbf{Religion of Life} and a \textbf{Superman} who, strange to say, almost always ends in associations of women and sub-men: from all of these forms a common meaning is revealed. 

It is the unraveling of the European soul, it is its self-release, its escape. 

Diverted by an insane drive for liberation, it evades the real, not for a super-real, but rather for a sub-real and a pre-real in which the sense of individuality is dissolved, and a cloudy, ecstatic coalescence with sub-human forces abolishes the law of pure action and clear vision. 

\begin{quotex}
With the void left by the loss of the Tradition in the Middle Ages, many false and counterfeit traditions rush in. No point to belabour them … you see them them as counterfeits or not. The reference to Ahriman is interesting, the result of his time with the Anthroposophists in the Ur and Krur groups.

As we would expect he opposes the path of ecstasy in favour of \textbf{pure action} and \textbf{clear vision}. The former follows the latter, lest we forget.

\end{quotex}
As little as that against which it reacts, such a spiritualism therefore constitutes a principle: it is not a sign of rebirth, but is instead — like that which already Asianized the Greco-Roman world in the Alexandrian period, and to which it so strangely resembles — a symptom of decline, an expression of rejection and resistance in the universal turmoil. 

\begin{quotex}
Evola later dedicated a book to this topic, Mask and Face of Contemporary Spiritualism, which included critiques of Theosophy, Anthroposophy, Psychoanalysis, and several other neo-spiritual movements. Here, he compares them to the various spiritual imports from Persia and India, following Alexander's conquests. Nevertheless, Evola seems to have later excluded other Asian influences such as Mithraism, early Buddhism, Taoism, Tantra, \emph{et al}. However, while he never advocated incorporating their cultural artifacts along with their metaphysical teachings. Contrast that with some contemporary neo-Traditionalists. 

\end{quotex}
Hence, sad omens loom over the Western world: since it is not a question of a circumstance of recent times, but rather the logical conclusion of the very principles on which this civilization has developed. In America — which is the most formidable of the newest barbarisms — is there perhaps not found the endpoint of the industrial course initiated by European Civilization? 

And in Bolshevism, which in a certain way constitutes a different form of the identical danger — does it not reveal perhaps the law in the materialistic social mask of that mysticism of community that, through the Christian subversion, swept away the individual, hierarchical, and imperial values of the Greco-Roman world? 

\begin{quotex}
This is old news by now, especially as Bolshevism as such is no longer a threat to Europe. Perhaps in his day, there seemed a viable Third Way, but is of little significance now. Even the opposition between capitalism and liberalism is a mirage, since they not just co-exist, but are actually in a symbiotic relationship, despite some rumblings from the left fringe and even less consequential clamoring on the various claimants to the real Right. The reference to the Christian subversion is unfortunate here, for several reasons.

\begin{enumerate}
\item First of all, he claims to be relying on principles rather than contingent events, a promise he breaks here. 
\item Secondly, it is misleading. Many assume that casting off Christianity will suffice to solve the problem. Yet, as Europe becomes more and more de-Christianized, the problem worsens. Rather than unveiling the ``inner pagan" in the European soul, is has brought about a spiritual vacuum filled by materialism, hedonism, or neo-spiritualism. 
\item It is historically untrue. The Greco part of the world continued its hierarchical and imperial values for another millennium. The Roman part managed to recover it and maintained it in one form or another for a thousand years. Rome was already in decline; a spiritually strong people don't go looking for novelty. The degeneration of castes had already occurred in the city-states. Rome was forced to make concessions to the mass of Plebes around the city. This suffices for now. 
\end{enumerate}
\end{quotex}
All this tells us how little there is to hope about the efficacy of a reaction. 

Still one generation — two at most — and every surviving possibility can be choked off, and nothing will stop this great dark mass that already runs down the slope: unless a sudden upheaval, a crisis that radically jolts the foundations of modern civilization comes to reestablish the equilibrium, whether through something, that through the eyes of most men for most of them will be the same as a catastrophe. 

Possessing this conviction, what task remains to the few that still resist? 

Not a direct action, but that more disconcerting action which can wield the mute and impassive presence of a stone guest. We must break the bridges, and with absolute adherence to primordial meanings and visions, those who will still act before the causes of the current civilization are established, constitute a pole, which, if it will not prevent this world of deviants from being what it is, will impede them, however, from asserting the nonexistence of every other horizon, to glorify itself, to establish by law itself to religion, to think that which is, should be what it is and that is good that it is. 

Hence, a fixed point, and from such a point, new relationships, new distances, new consciousnesses; from such consciousnesses, perhaps — in someone — the beginnings of liberating crises. 

\begin{quotex}
Well, we are two generations from the publication of Evola's essay. Are all possibilities choked off? Yet, don't we see more interest in Tradition, especially among young men? More translations are available, more commentaries on Eastern metaphysics, and so on. The World Cycle has its own program and is not as impatient as we are.

A catastrophe may occur but is not the only possibility as Guenon pointed out. The task is simple: \emph{absolute adherence to primordial meanings and visions}. Then act as a \textbf{stone guest}, a sojourner in the land of the lost, act as the unmoved mover, without attachment to results.

\end{quotex}
It is natural that many points in this regard be made precise and clarified: we will return to that task in our subsequent articles. Up till now we say that it is not a question of returning, because the reference is primarily to certain principles and certain interests, which, being above time, they have (in the words of Guenon) a permanent relevance.

\begin{quotex}
To be a man \textbf{above time}, adhere to principles whose relevancy and applicability are never in doubt. 

\end{quotex}
Having lost the sense of this relevance, having been dissolved in the myth of a pure flowing, of a pure stretching that pushes further and further its very goal of a process increasingly powerless to reach mastery, this is one of the characteristics of the world against which we anti-moderns set ourselves. 

Hence, a clean boundary that separates two epochs, not in an historical sense, but rather in an ideal sense: and we might call one the \textbf{Traditional}, the other \textbf{anti-traditional}. 

The first point is that relating to the spirit of the beginning, cancels every difference of the common opposition to the other. Then, we would like especially to mention the symbol closest to us Westerners: the symbol of \textbf{action}, restored to its full and traditional significance, of which the equivocal defense of the West of today could have a shapeless premonition. 

\begin{quotex}
To be a man \textbf{against time}, act on Traditional principles.

\end{quotex}
But not before the fixed point is established; that the sense of the distance is precise, so that the modality and nature of the processes appears, which sustain and foment the perversion of the European soul.



\flrightit{Posted on 2012-01-30 by Aeneas }

\begin{center}* * *\end{center}

\begin{footnotesize}\begin{sffamily}



\texttt{Avery on 2012-01-31 at 01:58 said: }

``The so-called crisis is entwined in the very structure of the modern world. It is not a question of fixing a limited problem, winning an election, or purging the West of some group or form of thought… We know things are awry, so what is the point of indulging in criticisms of this injustice or that stupid policy?"

A spot-on point that people know intuitively, and are being taught to forget by the media. Politics is one of the blasts that a culture unleashes into the world, like the corona of the Sun. It is an epiphenomenon and attacking who sits on the throne at the moment, on the basis of your cultural reaction to him, is the farthest thing from orienting yourself as an actor within your culture at large.


\hfill

\texttt{Avery on 2012-01-31 at 02:04 said: }

As for Evola's approach to Christendom, he strikes me as the epitome of European alienation. With alienation comes spiritual destruction, so Evola looks for a more ancient and mysterious spirit. His firm resolution to embrace both tradition and Europe without also embracing Christianity is the struggle behind much of his work.


\hfill

\texttt{Boreas on 2012-01-31 at 04:43 said: }

Well said, Avery. Evola was right and had superior knowledge on many things, but he forgot or did not fully understand – maybe because of his ``personal equation" which led to his errors also (manipura again) – the three things Christ brought to the world with the fulfillment of his own Great Work: Unity, Love (ultimate compassion and loving-wisdom) and the embodiment of the Spirit from the highest realms into the deepest recesses of matter.

Recommended reading: Matt. 12, the whole chapter


\hfill

\texttt{Cologero on 2012-02-01 at 22:42 said: }

``And why should it be that among the greatest mystics only the women, Saint Ignatius Loyola excepted, have been visionaries? Think about Catherine of Siena, Marie of the Incarnation, Mechtilde of Hackborn, Elizabeth of Schoenau, Julian of Norwich and others. For the most part, like oceans or tidal pools, women keep their vital secrets. \flright{\textsc{Evan S Connell, Points for a Compass Rose}}


\hfill

\texttt{Charlotte Cowell on 2012-02-02 at 16:26 said: }

yes, it is tricky, especially when you don't know someone, hence my checking in about possible upsets (I can be quite blunt, mainly I make myself laugh and not necessarily anyone else!). But regarding that comment I was not actually `worried' but it would be very complicated to convey why I was feeling beyond confirming I was very amused! However I can be a nervous type of person at times – highly strung you might say – so it is not surprising if you wondered if I really WAS worried!!

@Boreas, I went through a heavy metal phase when I was about 13 and shall never, ever return. However, I do think Led Zeppelin are an amazing band, one of the best ever, even though they're a bunch of lefthanders, well Jimmy Page at least. There's an extremely readable biography by Peter Levanda called Stairrway to Heaven, covering the whole band, which I really enjoyed. Levanda has also written equally gripping books on Hitler's Occult War – Unholy Alliance. Other than Led Zep a bit of Deep Purple is fine, same for ACDC, I'm not entirely writing off the genre, but I really prefer……D I S C O (the Studio 54 and Paradise Garage kind [D83D?][DE42?] ), world music, reggae, house, chillout stuff and things of that ilk. Some people call it `girly music', although the best DJs are usually men 

Point 2 you make I agree is valid in terms of catharsis, but there is also a danger of dark music or whatever invoking dark forces, as is the intention of the artist in many cases. In the same way you hear atonal rap music now, which is truly repulsive in most cases – I can't see this having any kind of good effect on listeners and the same goes for the worst kind of death metal, which can really be very sinister.


\hfill

\texttt{Charlotte Cowell on 2012-02-02 at 17:57 said: }

Boreas, with respect to something we were talking about on another thread, I read someone post this on facebook that is relevant to the points you were making: 

The gospel teaches that Yohanan the Baptist lived in the wilderness and that he ate ``locusts and wild honey." According to the Kabbalah, the ``locusts" represent admixed and dark forces, and the ability of the prophet to devour and transmute them, and the ``wild honey" represents divine forces, and the ability of the prophet to receive and integrate them, and embody them. In the Yeshua Messiah we are also called to this ``diet," empowered in the very same spiritual work so long as we live in this wilderness. Peace be with you!


\hfill


\hfill

\texttt{uberman on 2012-02-01 at 19:01 said: }

This series is a very welcome, perhaps belated diction and commentary on our historical predicament, penned form the heights. Like a breath of cool air in the encroaching desert, bringing with it a message of steadfastness and hope to an otherwise panic stricken cause. 

``If you can force your heart and nerve and sinew

To serve your turn long after they are gone,

And so hold on when there is nothing in you

Except the Will which says to them: `Hold on!'" – R.K


\end{sffamily}\end{footnotesize}
